\chapter{Introduction}\label{ch:intro}

\section{Motivation}

A poet's style is the part of a poem that survives a change of subject: the
words they reach for, how long their lines run, whether they rhyme and how,
which images return. Rewriting a poem so that its content stays and its
voice becomes another poet's is a hard test of whether a language model has
learned that layer of language rather than its surface. This is the task of
\emph{text style transfer}, here in its most demanding form: the ``style''
is not a register such as formal or informal, but one particular author.

Most work on style transfer and on language models in general is done in
English, where data is plentiful. Macedonian is a low-resource language: it
has few annotated corpora, few tools, and it is a small share of what
general-purpose models are trained on. Poetry makes this worse. A poet
writes a few hundred poems in a lifetime, often far fewer, so even a
complete collection of an author's work is a small dataset by the standards
of machine learning. The corpus this thesis is built on holds 1,239 poems
and song lyrics by 38 Macedonian authors; the best represented author has
182 poems, the next two 160 and 134, and half of the authors
have seventeen or fewer.

Large models can be prompted to write in a given style, but they are
expensive to run and are not under the user's control. Small models of a
few billion parameters can be trained and run on a single GPU, which makes
them the practical choice for a language community without large
computational resources. Whether they can learn a task as fine-grained as
an individual poet's voice, from so little data, is the question this
thesis sets out to answer.

\section{Research question}

The thesis asks:

\begin{quote}
\emph{Can a small language model, given only a source poem and the name of
a Macedonian author, rewrite the poem in that author's style, and does
reinforcement learning improve the result?}
\end{quote}

Two parts of the question need making precise. ``In that author's style''
is measured in two ways that do not depend on each other: automatically,
by stylometric features and by an author classifier that the training never
sees, and by human judgement. ``Reinforcement learning'' is Group Relative
Policy Optimization (GRPO) \pending{cite: Shao et al.\ 2024}, applied on top
of supervised fine-tuning with a reward built from the same stylometric
profile of each author.

\section{Approach and main finding}

The work reached its final form after several approaches that did not work.
Replacing selected words of the source poem with words typical of the
target author, inflected to fit, preserves the poem but changes little of
its style. Asking small local models to fill such slots fails at
Macedonian itself. What did work was a large model, Gemma~4 31B
\pending{cite}, prompted with a detailed scaffold: a summary of the source
poem, a prose description of the target author's style, a vocabulary
palette and example lines, and explicit structural targets. That model
became the \emph{teacher}. Its outputs, about 115 per target author for ten
target authors, are the training data for two \emph{students} of similar
size, Qwen3-4B \pending{cite} and Gemma~4 E4B, which see only the source
poem and the author's name. Each student is fine-tuned with LoRA
\pending{cite: Hu et al.\ 2022} on the teacher's poems and then trained
further with GRPO.

The main finding is negative. The students learn the task's form: they
write Macedonian, they write poems rather than prose, they change the
source rather than copy it. They do not learn the voices. An independent
author classifier recognises the intended author in 16.7\% of the
fine-tuned Qwen student's poems on the held-out test pairs, against 66.3\%
for real, unseen poems by the same authors; after GRPO the figure falls to
12.4\%, even though the reward it optimises improves.
\pending{Gemma student and teacher numbers on test; human ratings.} The
outputs that are identified belong almost entirely to the three authors
with the most poems in the corpus. The thesis argues that the limit is the
data rather than the method: a stronger student, a better-tuned reward and
the teacher itself all run into the same small, uneven corpus.

\section{Contributions}

\begin{itemize}
  \item A cleaned and part-of-speech-tagged corpus of Macedonian poetry,
        with poem-level identifiers, documented tagging corrections and a
        fixed train/validation/test split by author.
  \item Per-author stylometric profiles and a style-transfer reward built
        from them: part-of-speech distributions, lexical diversity, line
        and stanza structure, rhyme, characteristic vocabulary and an
        author classifier, with gates against copying the source and
        against degenerate output.
  \item A teacher--student pipeline that distils a large, heavily prompted
        model into two small students, followed by GRPO, run with matched
        settings for both students so that their results can be compared.
  \item An evaluation that separates real gains from reward gaming: the
        classifier used as a reward is never used to judge the outcome,
        a second classifier trained on different features is.
  \item A human evaluation of the teacher, both students and real poems on
        the same test pairs, with the real poems as a blind control.
\end{itemize}

\section{Structure of the thesis}

Chapter~\ref{ch:background} reviews text style transfer, stylometry,
knowledge distillation and reinforcement learning for language models.
Chapter~\ref{ch:data} describes the corpus and its preparation, and
Chapter~\ref{ch:stylometry} the stylometric profiles and classifiers the
rest of the work relies on. Chapter~\ref{ch:path} summarises the earlier
approaches and why they were abandoned. Chapter~\ref{ch:method} presents the
teacher, the students and the reward; Chapter~\ref{ch:eval} the
evaluation protocol. Chapter~\ref{ch:results} reports the results,
Chapter~\ref{ch:discussion} discusses why the students fall short, and
Chapter~\ref{ch:conclusion} concludes.
